\documentclass[11pt]{article}
\usepackage[margin=1in]{geometry}
\usepackage{amsmath,amssymb,amsthm}
\usepackage{graphicx}
\usepackage[hidelinks]{hyperref}

\newtheorem{theorem}{Theorem}

\newcommand{\norm}[1]{\left\lVert #1 \right\rVert}
\newcommand{\abs}[1]{\left\lvert #1 \right\rvert}
\newcommand{\esssup}{\operatorname*{ess\,sup}}

\title{The Sharp Constant Two in Talponen's ODE-Determined \(L^{p(\cdot)}\) Norms}
\author{Agent lane 15}
\date{June 29, 2026}

\begin{document}
\maketitle

\section*{Source and Target}

The source paper is Jarno Talponen, \emph{ODE to \(L^p\) norms},
arXiv:1402.0528.  The paper defines an ODE-determined variable exponent
norm by the terminal value \(\norm{f}_{p(\cdot)}=\varphi_f(1)\), where
\(\varphi_f\) is the stabilized Caratheodory solution of
\[
  \varphi(0)=0^+,\qquad
  \varphi'(t)=\frac{\abs{f(t)}^{p(t)}}{p(t)}\varphi(t)^{1-p(t)}
  \quad\text{for a.e. }t\in[0,1].
\]
After giving examples showing that the constant \(2\) would be best possible,
the source asks whether
\[
  \frac12\norm{f}_1\leq \norm{f}_{p(\cdot)}\leq 2\norm{f}_\infty
\]
always holds.

\begin{center}
\includegraphics[width=0.9\textwidth]{figures/open_problem_crop.png}
\end{center}

\section*{Idea of the Proof}

The upper estimate improves the source paper's \(e\norm{f}_\infty\) bound.
The source observes that after the profile reaches \(\norm{f}_\infty\), one
has \(\varphi'\leq \varphi\).  But at and above that level the stronger
estimate \(\varphi'\leq\norm{f}_\infty\) holds.  Thus the profile can gain at
most one more copy of \(\norm{f}_\infty\) during the remaining interval.

The lower estimate is most economical through Proposition 3.3 of the source
paper, which compares the ODE norm to the corresponding Nakano norm.  Young's
inequality gives \(\norm{f}_1\leq 2\) times the Nakano norm, and the Nakano
norm is bounded above by the ODE norm.

\begin{theorem}
Let \(p:[0,1]\to[1,\infty)\) be measurable and let \(f\in L^{p(\cdot)}\) in
Talponen's ODE-determined sense.  Then
\[
  \frac12\norm{f}_{L^1[0,1]}\leq \norm{f}_{p(\cdot)}
  \leq 2\norm{f}_{L^\infty[0,1]},
\]
with the usual interpretation that the upper inequality is vacuous if
\(\norm{f}_\infty=+\infty\).  The constant \(2\) is optimal on both sides.
\end{theorem}

\begin{proof}
The case \(f=0\) is trivial, so assume \(f\neq0\).  Write
\(\lambda=\norm{f}_{p(\cdot)}=\varphi_f(1)\).

First prove the upper bound.  Put \(M=\norm{f}_\infty\).  If \(M=+\infty\)
there is nothing to prove, and if \(M=0\) then \(f=0\), already excluded.
Assume \(0<M<\infty\).  If \(\varphi_f(1)\leq M\), then
\(\varphi_f(1)\leq 2M\).  Otherwise, by continuity and monotonicity of
\(\varphi_f\), choose \(t_0\in[0,1]\) such that \(\varphi_f(t_0)=M\).  For
almost every \(t\in[t_0,1]\) we have \(\varphi_f(t)\geq M\) and
\(\abs{f(t)}\leq M\), hence
\[
  \varphi_f'(t)
  =
  \frac{\abs{f(t)}^{p(t)}}{p(t)}\varphi_f(t)^{1-p(t)}
  \leq
  \frac{M^{p(t)}}{p(t)}M^{1-p(t)}
  =
  \frac{M}{p(t)}
  \leq M.
\]
Integrating over \([t_0,1]\) gives
\[
  \norm{f}_{p(\cdot)}=\varphi_f(1)
  \leq M+M(1-t_0)\leq 2M.
\]

Now prove the lower bound.  Let
\[
  |||f|||_{p(\cdot)}
  =
  \inf\left\{\mu>0:
    \int_0^1 \frac1{p(t)}
      \left(\frac{\abs{f(t)}}{\mu}\right)^{p(t)}\,dt
    \leq 1
  \right\}
\]
be the Nakano norm used in Proposition 3.3 of the source paper.  That
proposition states that
\[
  |||f|||_{p(\cdot)}\leq \norm{f}_{p(\cdot)}.
\]
Let \(\mu>|||f|||_{p(\cdot)}\) and set \(u(t)=\abs{f(t)}/\mu\).  By the
definition of the infimum, after decreasing \(\mu\) slightly if necessary, we
may assume
\[
  \int_0^1 \frac{u(t)^{p(t)}}{p(t)}\,dt\leq 1.
\]
For every \(q\geq1\) and \(s\geq0\),
\[
  s\leq \frac{s^q}{q}+1.
\]
Indeed, this is immediate for \(q=1\), and for \(q>1\) it follows from
Young's inequality \(s\leq s^q/q+1/q'\leq s^q/q+1\).  Applying this pointwise
with \(q=p(t)\) and \(s=u(t)\), and using that the measure of \([0,1]\) is
one, we get
\[
  \frac{\norm{f}_1}{\mu}
  =
  \int_0^1 u(t)\,dt
  \leq
  \int_0^1 \frac{u(t)^{p(t)}}{p(t)}\,dt+1
  \leq 2.
\]
Letting \(\mu\downarrow |||f|||_{p(\cdot)}\) yields
\(\norm{f}_1\leq 2|||f|||_{p(\cdot)}\leq 2\norm{f}_{p(\cdot)}\), which is the
desired lower bound.

The source paper's two examples immediately before the question show
sharpness.  With \(p\) equal to \(p_1\) on \([0,t_0)\) and \(1\) on
\([t_0,1]\), letting \(p_1\to\infty\) and \(t_0\to0\) gives
\(\norm{\mathbf1}_{p(\cdot)}\to2=2\norm{\mathbf1}_\infty\).  The altered
example \(f_{t_0}=1/t_0\) on \([0,t_0)\) and \(f_{t_0}=1\) on \([t_0,1]\),
with first exponent \(1\), second exponent tending to infinity, and
\(t_0\to0\), gives \(\norm{f_{t_0}}_{p(\cdot)}\to1\) while
\(\norm{f_{t_0}}_1\to2\).  Hence no constant smaller than \(2\) can replace
the right-hand constant, and no constant larger than \(1/2\) can replace the
left-hand constant.
\end{proof}

\section*{Verification and Novelty Notes}

The proof is analytic.  No numerical verification is used.  The only source
paper input beyond the definition of the ODE norm is Proposition 3.3, the
comparison \(|||f|||_{p(\cdot)}\leq\norm{f}_{p(\cdot)}\).

Local duplicate checks were performed on June 29, 2026 in
\texttt{registry\_index.tsv}, \texttt{solutions/index.tsv},
\texttt{attempts/index.tsv}, and \texttt{proof\_gaps/index.tsv} for arXiv id
1402.0528 and the phrases ``ODE'', ``variable exponent'', ``constant 2'',
``Nakano'', and variants of the displayed inequality.  No existing packet for
this target was found.  A bounded web search for the arXiv id, the paper
title, Talponen with ``constant 2'', and exact text from the displayed
inequality did not identify a later answer.  This is a bounded-search novelty
check, not a complete bibliographic audit.

\section*{Limitations and Human Review Recommendation}

The proof answers the exact constant-two problem stated in the source paper.
The reviewer should check only routine convention points: the interpretation
when \(\norm{f}_\infty=+\infty\), the use of Proposition 3.3 for the source
paper's ODE-determined domain, and the endpoint \(p(t)=1\) in the elementary
Young inequality step.  I recommend human review as a likely valid full
solution.

\begin{thebibliography}{9}
\bibitem{Talponen}
J. Talponen, \emph{ODE to \(L^p\) norms}, arXiv:1402.0528.
\end{thebibliography}

\end{document}
